\section{Superseded: verification by hand}\label{sec:hand}

Before \texttt{tools/cpp/goldbach.cpp} existed, the verified range was
$n\leq1000$, obtained by the argument below.  Theorem~\ref{verify:thm}
supersedes it completely.  It is kept because the method, not the range, is
what may be worth reusing: it needs no sieve, and it generalises to any
sufficiently dense set of summands.

\begin{prop}\label{hand:prop}
Every even $n$ with $4\leq n\leq1000$ is a \term{goldbach-number}.
\end{prop}

\begin{proof}[Proof as originally given]
Let $Q=\{3,5,7,11,13,17,19,23,29,31\}$.  For each even $n$ in range one checks
that some $q\in Q$ has $n-q$ prime; the verification is a table of $499$ rows,
which was checked by hand and is not reproduced.  The choice of $Q$ is not
canonical --- any set of small primes whose reflections cover the range will
do --- and it is exactly this arbitrariness that the sieve argument removes.
\end{proof}

\emph{Why it was replaced.}  The table does not scale, and its correctness
rests on a reader's patience rather than on anything checkable.  The
replacement is stronger in three ways: the range is $10^{5}$ times larger, the
witness bound of Theorem~\ref{verify:thm} is a genuine statistic rather than an
artefact of the chosen $Q$, and the computation is reproducible by one command.
That is the standard a computational claim in this project is held to.
